\subsection{失败了一整天之后，应该如何恢复？}

如果这一天已经失败了，第一业务目标不是把整个自我修好。
真正要做的，是防止这个失败状态明天继续复制自己。
恢复最有效的方式，是把这次失败当作一个状态问题来处理：识别当前状态，降低进入下一个可用目标状态的障碍，并只花足够的意志能量，让自己在下一次决策窗口里重新动起来。

\begin{enumerate}[nosep]
  \item 用操作语言命名这次失败。你错过了哪个目标状态，最后是哪个当前状态赢了？
  \item 找出让这次失败持续下去的主障碍。是奖赏俘获、模糊性、环境、疲劳，还是单纯错失了时机？
  \item 为下一步选一个更小的目标状态。不要试图一次性把整套计划都抢救回来，而要选那个能最低成本恢复接触的动作。
  \item 在睡前或停下前去掉一个摩擦源。把文档、衣物、工具或提醒放到位，让下一次决策窗口以更低阻力打开。
  \item 写下明天的第一个可见动作。这个动作要小到即使意志能量很低也能活下来。
  \item 标记下一次检查点。恢复不是“感觉好一点”就算完成，而是“下一次转变开始得更容易”才算完成。
\end{enumerate}

\begin{remark}
失败的一天看上去像道德崩塌，但从机制上说，它通常只是惯性模式暂时变强了。
最佳回应不是对过去的一天做一篇长长的忏悔，而是用一个短干预，让下一个目标状态比当前状态更容易被触发。
\end{remark}